\documentclass{article}
\usepackage[UTF8]{ctex}
\usepackage{amsmath,amssymb,amsfonts}
\usepackage{enumitem}
\usepackage{geometry}
\geometry{a4paper, margin=1in}
% === 错题本排版优化（不参与题目标识） ===
\setlength{\parindent}{0pt}
\setlength{\parskip}{0.35em plus 0.10em minus 0.08em}
\setlength{\emergencystretch}{3em}
\setlist{itemsep=0.25em, topsep=0.40em, parsep=0pt}
\allowdisplaybreaks[2]
% === 错题本排版优化结束 ===
\begin{document}
\section*{西南财经大学2024年数学分析真题}
\section*{173．西南财经大学2024年数学分析试卷}

\subsection{题目1}
1．求极限 \(\lim _{x \rightarrow \frac{\pi}{4}} \frac{\sec ^{2} x-2 \tan x}{1+\cos 4 x}\) 。

\subsection{题目2}
2．设函数 \(f(x)\) 在区间 \([0,1]\) 上二阶可导，且 \(f(0)=0, f(1)=1\) ．

\begin{itemize}
  \item[（1）] 证明：存在 \(x_{0} \in(0,1)\) ，使得 \(f^{\prime}\left(x_{0}\right)=1\) ．
  \item[（2）] 证明：存在 \(\xi \in(0,1)\) ，使得 \(\xi f^{\prime \prime}(\xi)+(1+\xi) f^{\prime}(\xi)=1+\xi\) 。
\end{itemize}

\subsection{题目3}
3．设函数

\[
f(x, y)= \begin{cases}x-y+\frac{x y^{3}}{x^{2}+y^{4}}, & (x, y) \neq(0,0) \\ 0, & (x, y)=(0,0)\end{cases}
\]

讨论 \(f(x, y)\) 在原点 \((0,0)\) 处是否连续，偏导数是否存在，是否可微．

\subsection{题目4}
4．设 \(z=z(x, y)\) 具有二阶连续偏导数，请计算 \(\frac{\partial^{2} z}{\partial x^{2}}+\frac{\partial^{2} z}{\partial y^{2}}\) 在坐标变换

\[
\left\{\begin{array}{l}
u=x^{2}-y^{2} \\
v=2 x y
\end{array}\right.
\]

下的表达式．

\subsection{题目5}
5．求 \(f(x, y)=2 x^{2}+y^{2}-8 x-2 y+9\) 在区域 \(D=\left\{(x, y) \mid 2 x^{2}+y^{2} \leq 1\right\}\) 上的最大值和最小值．

\subsection{题目6}
6．计算二重积分 \(I=\iint_{D}|x-y| \mathrm{d} x \mathrm{~d} y\) ，其中 \(D=\left\{(x, y) \mid x^{2}+y^{2} \leq a x\right\}, a>0\) ．

\subsection{题目7}
7．求下列级数的和 \(\sum_{n=2}^{\infty} \frac{1}{2^{n}\left(n^{2}-1\right)}\) ．

\subsection{题目8}
8．判断反常积分 \(\int_{0}^{+\infty} \frac{x^{p-1}}{1+x} \mathrm{~d} x\) 的敛散性．
\end{document}
